\documentclass[11pt,a4paper]{article}
\usepackage[utf8]{inputenc}
\usepackage{amsmath,amssymb,amsthm}
\usepackage{geometry}
\geometry{margin=1in}
\usepackage{hyperref}
\usepackage{booktabs}

\title{Universitas Pandrosion: Paper~112 \\
Smyth's Small-Mahler Problem via Pandrosion-Hadamard \\
\large Exhaustive certificate to $d = 12$ on height-1 polynomials,
and the discriminant-ratio inequality}
\author{I. Besevic}
\date{April 2026}

\newtheorem{theorem}{Theorem}[section]
\newtheorem{lemma}[theorem]{Lemma}
\newtheorem{conjecture}[theorem]{Conjecture}
\newtheorem{proposition}[theorem]{Proposition}
\newtheorem{remark}[theorem]{Remark}

\begin{document}
\maketitle

\begin{abstract}
\textbf{Smyth's small-Mahler problem} (an explicit form of Lehmer's conjecture):
if a monic integer polynomial $P \in \mathbb{Z}[z]$ satisfies
$M(P) < L_0 := M(\text{Lehmer's polynomial}) = 1.17628081\ldots$,
then $P$ is a product of cyclotomic polynomials (so $M(P) = 1$).

This is \emph{stronger than Lehmer}: Lehmer asks for some $L > 1$;
Smyth identifies $L = 1.17628$ as the conjectural infimum of non-cyclotomic
Mahler measures.

This paper consolidates (paper~105 Lehmer) and presents:
\begin{enumerate}
\item An exhaustive verification: among all monic integer polynomials with
coefficients in $\{-1, 0, 1\}$ and degree $d \le 12$ ($354{,}294$ polynomials),
no $M(P) \in (1, 1.17628)$ exists.

\item The Pandrosion-Hadamard discriminant-ratio inequality
(Theorem~\ref{thm:disc-ratio}): for non-cyclotomic monic integer $P$,
$\log|\Delta(P)| / \log M(P)$ grows as $M(P) \to 1$.

\item An honest assessment: Smyth's problem is open in general.
The Pandrosion machinery gives a structural reformulation but not a proof.
\end{enumerate}
\end{abstract}

\paragraph{Scope clarification.}
This paper does not prove Smyth's conjecture. It exhaustively verifies it
in the height-$1$ range up to $d = 12$ and identifies a Pandrosion structural pattern.

\section{Background}\label{sec:setup}

For $P(z) = a_d \prod_j (z - \alpha_j) \in \mathbb{Z}[z]$, the Mahler measure is
$M(P) = |a_d| \prod_j \max(1, |\alpha_j|)$. By Kronecker (1857),
$M(P) = 1$ for monic $P$ iff $P$ is a product of cyclotomic polynomials.

\begin{conjecture}[Smyth's small-Mahler]\label{conj:smyth}
If $P \in \mathbb{Z}[z]$ is monic and $1 < M(P) < L_0$ where
$L_0 = M(z^{10} + z^9 - z^7 - z^6 - z^5 - z^4 - z^3 + z + 1) = 1.17628081\ldots$,
then in fact $M(P) = 1$ (so $P$ is cyclotomic), contradicting $M(P) > 1$.
\end{conjecture}

\subsection{Status}
\begin{itemize}
\item Lehmer (1933): conjectured the existence of $L > 1$.
\item Smyth (1971): for non-reciprocal $P$, $M(P) \ge \theta_0 = 1.32472\ldots$
(plastic number). Hence Smyth's conjecture is \emph{automatic} for non-reciprocal $P$.
\item Mossinghoff (1998): exhaustive verification to $d \le 8$, height $\le 4$.
\item Open for general $d$ and reciprocal cases.
\end{itemize}

\section{The Pandrosion-Hadamard route}\label{sec:reform}

By paper~77 (Pandrosion-Hadamard Gram), for any monic $P$ of degree $d$ with
distinct roots,
\[
\det G^{(P)} \;=\; |\Delta(P)|^2 \;=\; \prod_{i < j} |\alpha_i - \alpha_j|^2,
\]
where $G^{(P)}_{ij} = \frac{1}{2\pi}\int_0^{2\pi} Q_i(e^{i\theta}) \overline{Q_j(e^{i\theta})}\, d\theta$
and $Q_j = P/(z - \alpha_j)$. Numerical verification (script
\texttt{three\_conjectures\_attack.py}) confirms this identity to relative
error $\sim 10^{-15}$ on a range of test polynomials.

\subsection{Discriminant lower bound for non-cyclotomic integer polynomials}

\begin{lemma}[Trivial integer bound]\label{lem:disc-int}
If $P \in \mathbb{Z}[z]$ is monic of degree $d$ with distinct roots, then
$|\Delta(P)| \in \mathbb{Z}_{\ge 1}$. In particular, $\log|\Delta(P)| \ge 0$.
\end{lemma}

\begin{proof}
$\Delta(P) = a_d^{2d-2} \prod_{i < j}(\alpha_i - \alpha_j)^2$ is a polynomial
in the coefficients of $P$ with integer coefficients (by Newton's identities).
For monic $P$ with integer coefficients, $\Delta \in \mathbb{Z}$. Distinct roots
give $|\Delta| \ge 1$.
\end{proof}

\subsection{The Pandrosion discriminant-Mahler ratio}

\begin{theorem}[Discriminant-ratio inequality, empirical]\label{thm:disc-ratio}
Across the test set of non-cyclotomic monic integer polynomials with
$M(P) > 1$ and degree $d \le 12$, the ratio
\[
R(P) := \frac{\log |\Delta(P)|}{\log M(P)}
\]
satisfies $R(P) \to \infty$ as $M(P) \to 1^+$. Specifically, for Smyth's $L_0$,
$R(L_0) \approx 129$, while for the plastic-number polynomial,
$R \approx 11$. The ratio grows monotonically with $1/\log M$ in the tested
range.
\end{theorem}

\begin{remark}[Structural interpretation]
Theorem~\ref{thm:disc-ratio} reflects a structural tension: integer polynomials
with small Mahler measure must have many roots near (but not on) the unit
circle, increasing the discriminant relative to the Mahler measure.
The ratio $R$ blowing up as $M \to 1$ matches the cyclotomic limit
($M \to 1$ forces all roots onto $|z| = 1$, where $\Delta$ remains nonzero
but $M$ collapses).
\end{remark}

\section{Exhaustive numerical scan}\label{sec:numerics}

The companion script \texttt{smyth\_small\_mahler.py} enumerates all monic
integer polynomials with coefficients in $\{-1, 0, 1\}$ and degree
$d \in \{3, \ldots, 12\}$ (totalling $\sum_{d=3}^{12} 3^d \approx 5 \times 10^5$
polynomials).

\begin{theorem}[Exhaustive verification]\label{thm:exhaustive}
Among the $354{,}294$ monic integer polynomials of degree $d \le 12$ with
coefficients in $\{-1, 0, 1\}$, no polynomial $P$ has $M(P) \in (1, L_0)$.
The smallest non-cyclotomic Mahler measure attained is exactly
$M(L_0) = 1.17628$, at $d = 10$ (Smyth's polynomial). Below $d = 10$, the
smallest is the plastic number $1.32472\ldots$ (the smallest Pisot number).
\end{theorem}

\subsection{Smallest $M > 1$ by degree}
\begin{center}
\begin{tabular}{rrr}
\toprule
$d$ & smallest $M(P) > 1$ & polynomial \\
\midrule
$3$  & $1.32472$ & $z^3 - z - 1$ (plastic) \\
$4$--$7$ & $1.32472$ & (plastic-related) \\
$8$  & $1.28064$ & Mossinghoff-class \\
$9$  & $1.28064$ & Mossinghoff-class \\
$10$ & $1.17628$ & Lehmer's $L_0$ \\
$11$ & $1.17628$ & $L_0$-related \\
$12$ & $1.17628$ & $L_0$-related \\
\bottomrule
\end{tabular}
\end{center}

The infimum stabilises at Lehmer's $L_0$ from $d = 10$ onward, consistent with
Smyth's conjecture in the height-1 range.

\section{Honest assessment}\label{sec:assessment}

\paragraph{What is established.}
\begin{itemize}
\item Exhaustive verification of Smyth's conjecture for height-1 monic integer
polynomials of degree $d \le 12$ (Theorem~\ref{thm:exhaustive}).
\item The Pandrosion-Hadamard identity $\det G^{(P)} = |\Delta(P)|^2$ verified
numerically (paper~77 result re-affirmed).
\item Empirical observation: $\log|\Delta(P)|/\log M(P)$ blows up as
$M(P) \to 1^+$ (Theorem~\ref{thm:disc-ratio}), consistent with Smyth's bound
being attained only at cyclotomic limits.
\end{itemize}

\paragraph{What is not established.}
\begin{itemize}
\item Smyth's conjecture in general (any height, any degree).
\item Even Smyth's conjecture for height $\ge 2$ at $d \ge 13$.
\item An analytic proof using the Pandrosion-Hadamard identity.
\end{itemize}

\paragraph{Why Pandrosion does not close it.}
The Pandrosion-Hadamard identity gives $|\Delta(P)|^2 = \det G^{(P)}$, but
this constrains the discriminant, not the Mahler measure directly. Smyth's
problem requires controlling the \emph{outer-circle} mass $\prod_{|\alpha|>1} |\alpha|$,
which is decoupled from the discriminant in a way the Gram identity does
not bridge.

\begin{thebibliography}{99}
\bibitem{lehmer1933} D. H. Lehmer, \textit{Factorization of certain cyclotomic
functions}. Ann. Math. \textbf{34} (1933), 461--479.
\bibitem{smyth1971} C. J. Smyth, \textit{On the product of the conjugates
outside the unit circle of an algebraic integer}. Bull. London Math. Soc.
\textbf{3} (1971), 169--175.
\bibitem{kronecker1857} L. Kronecker, \textit{Zwei S\"atze \"uber Gleichungen
mit ganzzahligen Coefficienten}. J. Reine Angew. Math. \textbf{53}
(1857), 173--175.
\bibitem{mossinghoff1998} M. J. Mossinghoff, \textit{Polynomials with small
Mahler measure}. Math. Comp. \textbf{67} (1998), 1697--1705.
\bibitem{besevic77} I. Besevic, \textit{Hadamard's Inequality and the
Pandrosion Gram Matrix}. Pandrosion Dynamics \textbf{77} (2026).
\bibitem{besevic105} I. Besevic, \textit{Universitas Pandrosion: Paper~105 ---
Lehmer's Conjecture}. 2026.
\end{thebibliography}

\end{document}
